% GENERATED FILE. Edit canonical lessons, then run npm run generate:books.
% canonical-source-hash: fnv1a64:a7c249791c5c0386
% canonical-lessons: GU-C532-posavu, GU-C532-gaju, GU-C532-nabhvu, GU-C532-anukul, GU-C532-rajinamu, GU-R532-review-525-532

\chapter{Can We Manage It?}
\label{ch:gu-can-we-manage-it}
\hlchaptermodality{532}

\begin{chapteropening}
\textbf{By the end of this chapter:} \emph{I can say whether we can afford or manage something.}

Complete the last lesson of chapter 532: I can say whether we can afford or manage something.
\end{chapteropening}

\section[\texorpdfstring{\gu{પોસાવું} (posāvũ)}{પોસાવું (posāvũ)}]{\gu{પોસાવું} (posāvũ) — to be within one's means}

\label{lesson:GU-C532-posavu}



\begin{quote}
Before the new one: say the Gujarati for a change, then the Gujarati for arrangements.
\end{quote}

\subsection*{You'll want to know: \gu{પોસાવું}}
\textbf{\gu{પોસાવું}} — \emph{posāvũ} — \textquotedblleft{}to be within one's means\textquotedblright{}.

\begin{quote}
\gu{જો} \gu{અમને} \gu{નવી} \gu{ગાડી} \gu{પોસાશે} \gu{તો} \gu{અમે} \gu{તે} \gu{ખરીદીશું}.
\end{quote}

\begin{quote}
\emph{jo amne navī gāḍī posāśe to ame te kharīdīśũ.}
\end{quote}

\begin{quote}
If we can afford a new car, we will buy it.
\end{quote}

Past (it …): \gu{તે} \gu{ભાડું} \gu{તેને} \gu{પોસાયું}, \emph{te bhāḍũ tene posāyũ}.

\begin{cousinweb}
\textbf{\gu{પોસાવું}} (\emph{posāvũ}) is filed by CDIAL under Sanskrit \emph{pōṣáyati}, \textquotedblleft{}nourishes\textquotedblright{}, beside a Gujarati \emph{posvũ}, \textquotedblleft{}to nourish\textquotedblright{} (CDIAL 8410); what one can keep up and pay for is affordable. The person takes \textbf{\gu{ને}}: \textbf{\gu{મને} \gu{પોસાય} \gu{છે}}.
\end{cousinweb}

\begin{grammarlens}[title={What you've built}]
One more word for saying what happens if.
\end{grammarlens}

\subsection*{Your turn}
\begin{itemize}
\raggedright
  \item \emph{Say it:} \emph{posāvũ}
  \item \emph{Say it:} \emph{posāvũ}, once more
  \item \emph{Say it:} \emph{jo amne navī gāḍī posāśe to ame te kharīdīśũ.} — If we can afford a new car, we will buy it.
  \item \emph{Recall:} say the Gujarati for a change, then the Gujarati for arrangements, then say \emph{posāvũ} again
\end{itemize}

\begin{tcolorbox}[breakable,colback=teal!4,colframe=teal!35!black,title={Before you move on}]
What does \gu{પોસાવું} mean? (\textquotedblleft{}To be within one's means\textquotedblright{}.) Say it once more.
\end{tcolorbox}
\section[\texorpdfstring{\gu{ગજું} (gajũ)}{ગજું (gajũ)}]{\gu{ગજું} (gajũ) — one's capacity}

\label{lesson:GU-C532-gaju}



\begin{quote}
Before the new one: say the Gujarati for arrangements, then the Gujarati for to be within one's means.
\end{quote}

\subsection*{You'll want to know: \gu{ગજું}}
\textbf{\gu{ગજું}} — \emph{gajũ} — \textquotedblleft{}one's capacity\textquotedblright{}.

\begin{quote}
\gu{જો} \gu{મારું} \gu{ગજું} \gu{હોત} \gu{તો} \gu{હું} \gu{મોટું} \gu{ઘર} \gu{ખરીદત}.
\end{quote}

\begin{quote}
\emph{jo mārũ gajũ hot to hũ moṭũ ghar kharīdat.}
\end{quote}

\begin{quote}
If I had the means, I would buy a big house.
\end{quote}

A neuter noun.

\begin{cousinweb}
\textbf{\gu{ગજું}} (\emph{gajũ}) is what a person can manage or afford: \textbf{\gu{મારું} \gu{ગજું} \gu{નથી}}, it is beyond me. No dictionary we use records an older form, so its history is not settled.
\end{cousinweb}

\begin{grammarlens}[title={What you've built}]
One more word for saying what happens if.
\end{grammarlens}

\subsection*{Your turn}
\begin{itemize}
\raggedright
  \item \emph{Say it:} \emph{gajũ}
  \item \emph{Say it:} \emph{gajũ}, once more
  \item \emph{Say it:} \emph{jo mārũ gajũ hot to hũ moṭũ ghar kharīdat.} — If I had the means, I would buy a big house.
  \item \emph{Recall:} say the Gujarati for arrangements, then the Gujarati for to be within one's means, then say \emph{gajũ} again
\end{itemize}

\begin{tcolorbox}[breakable,colback=teal!4,colframe=teal!35!black,title={Before you move on}]
What does \gu{ગજું} mean? (\textquotedblleft{}One's capacity\textquotedblright{}.) Say it once more.
\end{tcolorbox}
\section[\texorpdfstring{\gu{નભવું} (nabhvũ)}{નભવું (nabhvũ)}]{\gu{નભવું} (nabhvũ) — to get by}

\label{lesson:GU-C532-nabhvu}



\begin{quote}
Before the new one: say the Gujarati for to be within one's means, then the Gujarati for one's capacity.
\end{quote}

\subsection*{You'll want to know: \gu{નભવું}}
\textbf{\gu{નભવું}} — \emph{nabhvũ} — \textquotedblleft{}to get by\textquotedblright{}.

\begin{quote}
\gu{જો} \gu{પગાર} \gu{ઓછો} \gu{હશે} \gu{તો} \gu{આપણે} \gu{ઓછા} \gu{પૈસામાં} \gu{નભીશું}.
\end{quote}

\begin{quote}
\emph{jo pagār ocho haśe to āpaṇe ochā paisāmā̃ nabhīśũ.}
\end{quote}

\begin{quote}
If the pay is low, we will get by on less money.
\end{quote}

Past (it …): \gu{તે} \gu{કુટુંબ} \gu{નભી} \gu{ગયું}, \emph{te kuṭumb nabhī gayũ}.

\begin{cousinweb}
\textbf{\gu{નભવું}} (\emph{nabhvũ}) is inherited from Sanskrit \emph{nírvahati}, \textquotedblleft{}carries out, carries through\textquotedblright{}, and CDIAL lists the Gujarati verb, \textquotedblleft{}to endure, continue\textquotedblright{} (CDIAL 7397).
\end{cousinweb}

\begin{grammarlens}[title={What you've built}]
One more word for saying what happens if.
\end{grammarlens}

\subsection*{Your turn}
\begin{itemize}
\raggedright
  \item \emph{Say it:} \emph{nabhvũ}
  \item \emph{Say it:} \emph{nabhvũ}, once more
  \item \emph{Say it:} \emph{jo pagār ocho haśe to āpaṇe ochā paisāmā̃ nabhīśũ.} — If the pay is low, we will get by on less money.
  \item \emph{Recall:} say the Gujarati for to be within one's means, then the Gujarati for one's capacity, then say \emph{nabhvũ} again
\end{itemize}

\begin{tcolorbox}[breakable,colback=teal!4,colframe=teal!35!black,title={Before you move on}]
What does \gu{નભવું} mean? (\textquotedblleft{}To get by\textquotedblright{}.) Say it once more.
\end{tcolorbox}
\section[\texorpdfstring{\gu{અનુકૂળ} (anukūḷ)}{અનુકૂળ (anukūḷ)}]{\gu{અનુકૂળ} (anukūḷ) — convenient}

\label{lesson:GU-C532-anukul}



\begin{quote}
Before the new one: say the Gujarati for one's capacity, then the Gujarati for to get by.
\end{quote}

\subsection*{You'll want to know: \gu{અનુકૂળ}}
\textbf{\gu{અનુકૂળ}} — \emph{anukūḷ} — \textquotedblleft{}convenient\textquotedblright{}.

\begin{quote}
\gu{જો} \gu{તમને} \gu{રવિવાર} \gu{અનુકૂળ} \gu{હશે} \gu{તો} \gu{આપણે} \gu{ત્યારે} \gu{મળીશું}.
\end{quote}

\begin{quote}
\emph{jo tamne ravivār anukūḷ haśe to āpaṇe tyāre maḷīśũ.}
\end{quote}

\begin{quote}
If Sunday suits you, we will meet then.
\end{quote}

\begin{cousinweb}
\textbf{\gu{અનુકૂળ}} (\emph{anukūḷ}) is a tatsama, taken from Sanskrit \emph{anukūla}, \textquotedblleft{}following the slope of the bank\textquotedblright{}, and so \textquotedblleft{}favourable, suitable\textquotedblright{} (Monier-Williams).
\end{cousinweb}

\begin{grammarlens}[title={What you've built}]
One more word for saying what happens if.
\end{grammarlens}

\subsection*{Your turn}
\begin{itemize}
\raggedright
  \item \emph{Say it:} \emph{anukūḷ}
  \item \emph{Say it:} \emph{anukūḷ}, once more
  \item \emph{Say it:} \emph{jo tamne ravivār anukūḷ haśe to āpaṇe tyāre maḷīśũ.} — If Sunday suits you, we will meet then.
  \item \emph{Recall:} say the Gujarati for one's capacity, then the Gujarati for to get by, then say \emph{anukūḷ} again
\end{itemize}

\begin{tcolorbox}[breakable,colback=teal!4,colframe=teal!35!black,title={Before you move on}]
What does \gu{અનુકૂળ} mean? (\textquotedblleft{}Convenient\textquotedblright{}.) Say it once more.
\end{tcolorbox}
\section[\texorpdfstring{\gu{રાજીનામું} (rājīnāmũ)}{રાજીનામું (rājīnāmũ)}]{\gu{રાજીનામું} (rājīnāmũ) — a resignation}

\label{lesson:GU-C532-rajinamu}



\begin{quote}
Before the new one: say the Gujarati for to get by, then the Gujarati for convenient.
\end{quote}

\subsection*{You'll want to know: \gu{રાજીનામું}}
\textbf{\gu{રાજીનામું}} — \emph{rājīnāmũ} — \textquotedblleft{}a resignation\textquotedblright{}.

\begin{quote}
\gu{જો} \gu{સાહેબ} \gu{મારી} \gu{વાત} \gu{નહીં} \gu{માને} \gu{તો} \gu{હું} \gu{રાજીનામું} \gu{આપીશ}.
\end{quote}

\begin{quote}
\emph{jo sāheb mārī vāt nahī̃ māne to hũ rājīnāmũ āpīś.}
\end{quote}

\begin{quote}
If the boss won't listen to me, I will resign.
\end{quote}

A neuter noun.

\begin{cousinweb}
\textbf{\gu{રાજીનામું}} (\emph{rājīnāmũ}) is usually traced to Persian \emph{rāzīnāma}, \textquotedblleft{}a letter of consent\textquotedblright{}; in Gujarati it is the letter that gives up a post, and one gives it: \textbf{\gu{રાજીનામું} \gu{આપવું}}.
\end{cousinweb}

\begin{grammarlens}[title={What you've built}]
One more word for saying what happens if.
\end{grammarlens}

\subsection*{Your turn}
\begin{itemize}
\raggedright
  \item \emph{Say it:} \emph{rājīnāmũ}
  \item \emph{Say it:} \emph{rājīnāmũ}, once more
  \item \emph{Say it:} \emph{jo sāheb mārī vāt nahī̃ māne to hũ rājīnāmũ āpīś.} — If the boss won't listen to me, I will resign.
  \item \emph{Recall:} say the Gujarati for to get by, then the Gujarati for convenient, then say \emph{rājīnāmũ} again
\end{itemize}

\begin{tcolorbox}[breakable,colback=teal!4,colframe=teal!35!black,title={Before you move on}]
What does \gu{રાજીનામું} mean? (\textquotedblleft{}A resignation\textquotedblright{}.) Say it once more.
\end{tcolorbox}
\section[(dialogue)]{Looking back: from Carving and Shaping to Can We Manage It?}

\label{lesson:GU-R532-review-525-532}



\begin{quote}
Say the Gujarati for a sculpture and to carve.
\end{quote}

\subsection*{Your turn}
Say these aloud:

\begin{itemize}
\raggedright
  \item a sculpture, to carve, a shape, to frame (a picture), an exhibition
  \item a musical note, a rhythm, a singer, a raga, to strike up
  \item a wrestling ground, a wrestler, a rival, to throw down, to tire out (an opponent)
  \item a list, a budget, an arrangement, a diary, to organise
  \item a stage (of a process), a sequence, an estimate, half-finished, final
  \item a time limit, urgency, extra time, to push back, at long last
  \item a hassle, unexpected, sudden, a change, arrangements
  \item to be within one's means, one's capacity, to get by, convenient, a resignation
\end{itemize}

\begin{tcolorbox}[breakable,colback=teal!4,colframe=teal!35!black,title={Before you move on}]
A sculpture? (\textbf{\gu{શિલ્પ}}, \emph{śilp}.) To carve? (\textbf{\gu{કોતરવું}}, \emph{kotarvũ}.) A shape? (\textbf{\gu{આકાર}}, \emph{ākār}.) To frame (a picture)? (\textbf{\gu{મઢવું}}, \emph{maḍhvũ}.) An exhibition? (\textbf{\gu{પ્રદર્શન}}, \emph{pradarśan}.) A musical note? (\textbf{\gu{સૂર}}, \emph{sūr}.) A rhythm? (\textbf{\gu{તાલ}}, \emph{tāl}.) A singer? (\textbf{\gu{ગાયક}}, \emph{gāyak}.) A raga? (\textbf{\gu{રાગ}}, \emph{rāg}.) To strike up? (\textbf{\gu{છેડવું}}, \emph{cheḍvũ}.) A wrestling ground? (\textbf{\gu{અખાડો}}, \emph{akhāḍo}.) A wrestler? (\textbf{\gu{પહેલવાન}}, \emph{pahelvān}.) A rival? (\textbf{\gu{હરીફ}}, \emph{harīph}.) To throw down? (\textbf{\gu{પછાડવું}}, \emph{pachāḍvũ}.) To tire out (an opponent)? (\textbf{\gu{હંફાવવું}}, \emph{hamphāvvũ}.) A list? (\textbf{\gu{યાદી}}, \emph{yādī}.) A budget? (\textbf{\gu{બજેટ}}, \emph{bajeṭ}.) An arrangement? (\textbf{\gu{વ્યવસ્થા}}, \emph{vyavasthā}.) A diary? (\textbf{\gu{ડાયરી}}, \emph{ḍāyrī}.) To organise? (\textbf{\gu{યોજવું}}, \emph{yojvũ}.) A stage (of a process)? (\textbf{\gu{તબક્કો}}, \emph{tabakko}.) A sequence? (\textbf{\gu{ક્રમ}}, \emph{kram}.) An estimate? (\textbf{\gu{અંદાજ}}, \emph{andāj}.) Half-finished? (\textbf{\gu{અધૂરું}}, \emph{adhūrũ}.) Final? (\textbf{\gu{અંતિમ}}, \emph{antim}.) A time limit? (\textbf{\gu{અવધિ}}, \emph{avdhi}.) Urgency? (\textbf{\gu{તાકીદ}}, \emph{tākīd}.) Extra time? (\textbf{\gu{મહેતલ}}, \emph{mahetal}.) To push back? (\textbf{\gu{ઠેલવું}}, \emph{ṭhelvũ}.) At long last? (\textbf{\gu{આખરે}}, \emph{ākhre}.) A hassle? (\textbf{\gu{ઝંઝટ}}, \emph{jhanjhaṭ}.) Unexpected? (\textbf{\gu{અણધાર્યું}}, \emph{aṇdhāryũ}.) Sudden? (\textbf{\gu{ઓચિંતું}}, \emph{ocintũ}.) A change? (\textbf{\gu{ફેરફાર}}, \emph{pherphār}.) Arrangements? (\textbf{\gu{બંદોબસ્ત}}, \emph{bandobast}.) To be within one's means? (\textbf{\gu{પોસાવું}}, \emph{posāvũ}.) One's capacity? (\textbf{\gu{ગજું}}, \emph{gajũ}.) To get by? (\textbf{\gu{નભવું}}, \emph{nabhvũ}.) Convenient? (\textbf{\gu{અનુકૂળ}}, \emph{anukūḷ}.) A resignation? (\textbf{\gu{રાજીનામું}}, \emph{rājīnāmũ}.)
\end{tcolorbox}
